
\medskip

\emph{Principaux domaines abordés}: Manipulation des vecteurs, des droites et des plans de l'espace. Orthogonalité et distances dans l'espace.
Représentations paramétriques et équations cartésiennes.

\bigskip

Dans un repère orthonormé \Oijk{} de l'espace,
on considère les points 

\[\text{A}(-3~;~1~;~3), \text{B}(2~;~2~;~3), \text{C}(1~;~7~;~-1), \text{D}(-4~;~6~;~-1)\: \text{et 
K}(-3~;~14~;~14).\]

%\medskip

\begin{enumerate}
\item 
	\begin{enumerate}
		\item %Calculer les coordonnées des vecteurs $\vect{\text{AB}},\: \vect{\text{OC}}$ et $\vect{\text{AD}}$.
On a $\vect{\text{AB}}\begin{pmatrix}5\\1\\0\end{pmatrix}, \: \vect{\text{DC}}\begin{pmatrix}5\\1\\0\end{pmatrix}$ et $\vect{\text{AD}}\begin{pmatrix}-1\\5\\-4\end{pmatrix}$.
		\item% Montrer que le quadrilatère ABCD est un rectangle.
D'après la question précédente, $\vectt{AB}=\vectt{DC}$ donc ABCD est un parallélogramme.

De plus le produit scalaire $\vect{\text{AB}} \cdot \vect{\text{AD}} = 5\times(-1) + 1\times 5 + 0\times(-4) = 0$.

Le parallélogramme ABCD ayant deux côtés consécutifs perpendiculaires est un rectangle.

		\item %Calculer l'aire du rectangle ABCD.
On a AB$^2 = 25 + 1 = 26$, d'où AB $= \sqrt{26}$ ; \\
de même AD$^2 =  1  + 25 + 16= 42$, d'où AD $= \sqrt{42}$.

L'aire du rectangle ABCD est égale à \\
$\text{AB} \times \text{AD} = \sqrt{26 }\times\sqrt{42} = \sqrt{26 \times 42} = \sqrt{2 \times 13 \times 2 \times 21} = 2\sqrt{13 \times 21} = 2\sqrt{273}$.
	\end{enumerate}
\item 
	\begin{enumerate}
		\item %Justifier que les points A, B et D définissent un plan.
D'après la question 1., les vecteurs $\vect{\text{AB}}$ et $\vect{\text{AD}}$ ne sont pas colinéaires, donc les points A, B et D ne sont pas alignés : ils définissent donc bien un plan.
		\item Soit le vecteur $\vect{n}(-2~;~10~;~13)$.% est un vecteur normal au plan (ABD).
		
$\vect{n} \cdot \vect{\text{AB}} = - 10 + 10 + 0 = 0$ et $\vect{n} \cdot \vect{\text{AD}} = 2 + 50 - 52 = 0$.

Conclusion : le vecteur $\vect{n}$, orthogonal à deux vecteurs non colinéaires du plan (ABD), est normal à ce plan.
		\item %En déduire une équation cartésienne du plan (ABD).
Le résultat précédent montre que le plan (ABD) a une équation de la forme\\
$-2x + 10y + 13z = d$, avec $d \in \R$.

Or par exemple $\text{B}(2~;~2~;~3) \in (\text{ABD})$ donc \\
$-2 \times 2  + 10 \times 2  + 13\times 3 = d \iff - 4 + 20 + 39 = d \iff d = 55$.

Donc le plan (ABD) a pour équation $ -2x + 10y + 13z = 55$.
	\end{enumerate}
	
\item
	\begin{enumerate}
		\item %Donner une représentation paramétrique de la droite $\Delta$ orthogonale au plan (ABD) et qui passe par le point K.
Si $\Delta$ est orthogonale au plan (ABD) elle a pour vecteur directeur le vecteur $\vect{n}$. Comme elle contient K, on a donc :

$M(x~,~y~;~z) \in \Delta \iff \vect{\text{K}M} = t\vect{n}$, \: $t \in \R$. 

Avec $\vect{\text{K}M}\begin{pmatrix}x - (-3)\\y - 14\\z - 14\end{pmatrix}$ ceci se traduit par le système :

$\left\{\begin{array}{l c l}
x + 3&=&- 2t\\
y - 14&=&10t\\
z - 14&=&13t
\end{array}\right., t \in \R \iff \left\{\begin{array}{l c r}
x &=&- 2t - 3\\
y &=&10t + 14\\
z &=&13t + 14
\end{array}\right., t \in \R $
		\item %Déterminer les coordonnées du point I, projeté orthogonal du point K sur le plan (ABD).
Si I est le projeté orthogonal de K sur le plan (ABD), le point I est un point de $\Delta$; comme c'est aussi un point de (ABD) ; ses coordonnées vérifient  le système

$\left\{\begin{array}{l c r}
x &=&- 2t - 3\\
y &=&10t + 14\\
z &=&13t + 14\\
- 2x + 10y + 13z &=& 55
\end{array}\right., t \in \R $.

En remplaçant dans la dernière équation $x,\:y$ et $z$ par leurs expressions en fonction de $t$, on obtient :

$- 2(- 2t - 3) + 10(10t + 14) + 13(13t + 14) = 55 \\
\iff 4t + 6 + 100t + 140 + 169t + 182 = 55 \iff 273t = - 273 \iff t = - 1$.

Les premières équations donnent alors\\
 $x = 2 - 3 = - 1$, \: $y = - 10 + 14 = 4$ et $z = - 13 + 14 = 1$. 

Le point I a donc pour coordonnées $(- 1~;~4~;~1)$.


		\item %Montrer que la hauteur de la pyramide KABCD de base ABCD et de sommet K vaut $\sqrt{273}$.
ABCD étant un rectangle le point D appartient au plan (ABD). Donc la hauteur  de la pyramide KABCD de base ABCD est [KI].

$\vect{\text{KI}}\begin{pmatrix}2\\-10\\-13\end{pmatrix}$ donne KI$^2 = 4 + 100 + 169 = 273$ et enfin KI $= \sqrt{273}$.
	\end{enumerate}
\item %Calculer le volume $V$ de la pyramide KABCD.

%On rappelle que le volume V d'une pyramide est donné par la formule:

%\[V= \dfrac13 \times \text{aire de la base} \times \text{hauteur}.\]
On a $V = \dfrac13 \times \mathcal{A}(\text{ABCD}) \times \text{KI} = \dfrac13 \times 2\sqrt{273} \times \sqrt{273} = \dfrac{2 \times 273}{3} = \dfrac{2 \times 3 \times 91}{3} = 182$.
\end{enumerate}


